\title{A Friendly Guide to Loudspeakers:\\
From Small Signal T/S Parameters to Breakup Physics and Beyond}
\date{\today}
\author{Ivan Rodionov}

\begin{document}
\maketitle


\newpage
\section*{Overview}
This document provides a stepwise approach to loudspeaker modeling with clear explanations on what is assumed, what is added, and how the T/S framework scales toward full-featured loudspeaker modeling. Throughout we adopt the Laplace-domain notation (complex frequency variable $s$) except where explicitly stated otherwise.

\begin{enumerate}
    \item \textbf{Thiele--Small foundation:} Start with the small-signal T/S model, relating input voltage $V_{\mathrm{in}}(s)$ to on-axis sound pressure $p(s)$ via lumped electrical and mechanical parameters.  
    \item \textbf{Explicit assumptions:} Identify linearity, rigid-piston behavior, infinite baffle conditions, and far-field observation as the baseline simplifications.  
    \item \textbf{Scaling to reality:} Extend the model to include voice-coil inductance, radiation loading, and directivity effects.  
    \item \textbf{Generalization:} Reformulate in a modular, simulation-ready form that can be applied to arbitrary geometries, finite baffles, and code implementations.  
    \item \textbf{Practical Code Implementations:} Implement the derivations into an easy to use form to extract TS++ parameters for a given driver and for a given parameter list produce performance figures.  
\end{enumerate}

\section*{Notation \& Assumptions}
\begin{itemize}
  \item Linear, time-invariant small-signal behavior unless otherwise noted.  
  \item Constant motor strength $B\ell$ (no Bl(x) nonlinearity) for small excursions.  
  \item Infinite, rigid baffle for "baffled piston" analytic formulas (half-space radiation).  
  \item Far-field (Fraunhofer) approximations use $r\gg a^2/\lambda$ unless near-field terms are explicitly kept.  
  \item Laplace domain: $s$ is used everywhere (if you want steady-state set $s\mapsto j\omega$).
\end{itemize}

\section{Symbol Table (select)}
\begin{center}
\begin{tabular}{llc}
\toprule
Symbol & Meaning & Units \\
\midrule
$V_{\mathrm{in}}(s)$ & input voltage & V \\
$I(s)$ & coil current & A \\
$Z_e(s)$ & electrical impedance (coil + electromechanical reflection) & $\Omega$ \\
$B\ell$ & force factor (force per unit current) & $\mathrm{N\,A^{-1}}$ \\
$S_d$ & piston radiating area & $\mathrm{m^2}$ \\
$a$ & equivalent piston radius ( $S_d=\pi a^2$ ) & m \\
$M_{\mathrm{ms}}$ & moving mass of driver (including cone/coil) & kg \\
$C_{\mathrm{ms}}$ & mechanical compliance & m/N \\
$R_{\mathrm{ms}}$ & mechanical losses (damping) & N s/m \\
$Z_{\mathrm{rad}}^{\mathrm{ac}}(s)$ & acoustic radiation impedance (pressure / velocity) & Pa\;s/m \\
$Z_{\mathrm{rad}}^{\mathrm{mech}}(s)$ & mechanical radiation impedance (force / velocity) & N\;s/m \\
$k$ & acoustic wavenumber (Laplace) $k=\dfrac{s}{c}$ & $\mathrm{m^{-1}}$ \\
$\rho_0$ & air density & kg/m$^3$ \\
$c$ & speed of sound & m/s \\
\bottomrule
\end{tabular}
\end{center}


\newpage
\section{Impedance Definitions}

\subsection*{A. Electrical Impedance}
Electrical impedance describes how the voice coil of the driver resists and reacts to an applied voltage across frequency. It is the combination of the coil’s inherent resistance and inductance, along with the reflected impedance of the driver’s moving system through the force factor (Bl). This reflection links the mechanical domain back into the electrical domain, meaning any change in the cone’s motion affects the measured electrical response. 
\[
Z_e(s) = R_e + sL_e + \frac{(B\ell)^2}{Z_m^{\text{tot}}(s)}
\]
where $Z_m^{\text{tot}}(s) = Z_m(s) + Z_{\mathrm{rad, mech}}(s)$ is the total mechanical impedance.

\subsection*{B. Mechanical (Motional) Impedance}
Mechanical impedance expresses how the driver’s moving assembly resists motion when a force is applied. It combines mechanical resistance Rms (losses in suspension and surround), mass Mms (inertia), and compliance Cms (spring-like suspension). This impedance governs how easily the cone can accelerate or decelerate in response to the input signal, and forms the bridge between the electrical input and acoustic output. 
\begin{align*}
Z_m(s) &= R_{\mathrm{ms}} + sM_{\mathrm{ms}} + \frac{1}{sC_{\mathrm{ms}}}
\end{align*}
\subsection*{C. Radiation Impedance (Mechanical and Acoustic Domain) from Rayleigh}
This section starts from the Rayleigh (Kirchhoff) surface integral and follows a single logical path to the practical radiation models used in loudspeaker work: general operator, monopole limit, baffled-piston exact form, small-\(ka\) series (showing the algebraic step that yields the LF resistance and added mass), far-field algebraic simplification, conversion to mechanical radiation impedance, and finally the numerical ring-BEM fallback. All expressions are given in the Laplace domain with \(k=s/c\).

\subsection{Rayleigh integral — general acoustic operator.}
For an arbitrary radiating surface \(S\) with surface-normal velocity \(v(\mathbf{r}',s)\) the Laplace-domain Rayleigh integral gives the pressure at a field point \(\mathbf{r}\):
\[
\boxed{\;
p(\mathbf{r},s)
=\frac{s\rho_0}{2\pi}\iint_{S} v(\mathbf{r}',s)\,\frac{e^{-k|\mathbf{r}-\mathbf{r}'|}}{|\mathbf{r}-\mathbf{r}'|}\,dS'
\;}
\qquad\bigl(k=\tfrac{s}{c}\bigr).
\]
Define the average surface velocity and radiator area
\[
v_{\mathrm{avg}}(s)=\frac{1}{S_d}\iint_S v(\mathbf{r}',s)\,dS',\qquad S_d=\iint_S dS'.
\]
Then the acoustic velocity-to-pressure operator at \(\mathbf{r}\) (an impedance-like operator) is
\[
\boxed{\,Z_{\mathrm{rad}}^{\mathrm{ac}}(s;\mathbf{r})
\;=\;\frac{p(\mathbf{r},s)}{v_{\mathrm{avg}}(s)}
\;=\;\frac{s\rho_0}{2\pi}\frac{1}{v_{\mathrm{avg}}(s)}\iint_S v(\mathbf{r}',s)\frac{e^{-k|\mathbf{r}-\mathbf{r}'|}}{|\mathbf{r}-\mathbf{r}'|}\,dS'.\,}
\]

\subsection{Mechanical radiation impedance (force/velocity).}
The mechanical radiation impedance (which enters the motional branch of the driver) is the total radiated force per (average) velocity:
\[
\boxed{\,Z_{\mathrm{rad}}^{\mathrm{mech}}(s)=\frac{F_{\mathrm{rad}}(s)}{v_{\mathrm{avg}}(s)}
=\frac{1}{v_{\mathrm{avg}}(s)}\iint_S p_{\mathrm{near}}(\mathbf{r}',s)\,dS'\,}
\]
(units N·s/m). Under the uniform-face-pressure approximation (rigid piston with nearly uniform pressure) the two operators are related simply by area:
\[
\boxed{\,Z_{\mathrm{rad}}^{\mathrm{mech}}(s)\approx S_d\,Z_{\mathrm{rad}}^{\mathrm{ac}}(s)\,}
\]
(i.e. force \(\approx\) pressure \(\times\) area). Be explicit which form you use when assembling \(Z_m^{\mathrm{tot}}(s)\).

\subsection{Monopole (spherical) limit — small radiator, far-field.}
If the radiator is compact relative to the wavelength the total volume velocity is
\[
Q(s)=\iint_S v(\mathbf{r}',s)\,dS' = S_d\,v_{\mathrm{avg}}(s),
\]
and the Laplace-domain far-field on-axis pressure of a point monopole at distance \(r\) is
\[
\boxed{\,p_{\mathrm{mono}}(r,s)=\rho_0\,s\,\frac{Q(s)}{4\pi r}\,e^{-s r/c}
=\rho_0\,s\,\frac{S_d}{4\pi r}\,v_{\mathrm{avg}}(s)\,e^{-s r/c}\,}
\]
so the monopole acoustic impedance (velocity-to-pressure) is
\[
\boxed{\,Z_{\mathrm{mono}}^{\mathrm{ac}}(s,r)=\rho_0\,s\,\frac{S_d}{4\pi r}\,e^{-s r/c}\,.}
\]
This is causal (contains \(e^{-s r/c}\)) and useful as the small-source algebraic approximation.

\subsection{Baffled circular piston — exact on-axis (Bessel/Hankel) form.}
For a rigid circular piston of radius \(a\) in an infinite rigid baffle the Rayleigh integral evaluates on-axis to a closed-form involving Bessel functions. With \(S_d=\pi a^2\) and \(k=s/c\):
\[
\boxed{\,Z_{\mathrm{piston}}^{\mathrm{ac}}(s)
= \rho_0 c S_d\Bigl(1 - \frac{J_1(2ka)}{ka}\Bigr)
\;+\; \rho_0 c S_d\Bigl(\frac{2}{\pi k a}Y_1(2ka)\Bigr)\,}
\]
This is the exact (on-axis) acoustic velocity-to-pressure operator as a complex function of \(s\). On the \(s=j\omega\) axis the first bracket produces the resistive contribution and the second bracket the reactive contribution.\\
\textbf{Numerical note:} evaluate the exact Bessel/Hankel expression carefully for small \(|ka|\). Although not singular, the representation contains terms with \(1/(ka)\) that cause cancellation and loss of precision as \(ka\to 0\). Use series expansions (Taylor) for \(|ka|\ll1\) or robust special-function libraries that support complex arguments.

\subsection{Small-\(ka\) (Taylor) expansion — algebraic derivation of LF resistance and added mass.}
We now show how the exact piston form reduces to the familiar low-frequency lumped expressions by algebraic expansion.

Start with the small-argument expansion for \(J_1(x)\) (with \(x=2ka\)):
\[
J_1(x)=\frac{x}{2}-\frac{x^3}{16}+O(x^5).
\]
Hence
\[
\frac{J_1(2ka)}{ka}=1-\frac{(ka)^2}{2}+O((ka)^4).
\]
Substitute this into the first bracket of the exact piston expression:
\[
\rho_0 c S_d\Bigl(1-\frac{J_1(2ka)}{ka}\Bigr)\approx \rho_0 c S_d\cdot\frac{(ka)^2}{2}.
\]
Thus the leading low-frequency radiation resistance is
\[
\boxed{\,R_{\mathrm{rad}}(s)\approx \rho_0 c S_d\;\frac{(ka)^2}{2}
=\rho_0 c S_d\;\frac{(s a / c)^2}{2}\,.}
\]

The leading reactive (near-field / added-mass) term comes from the matched expansion of the second (Hankel/Y\(_1\)) term. Carrying out this matched-series expansion (standard result) yields
\[
\boxed{\,X_{\mathrm{rad}}(s)\approx \rho_0 c S_d\;\frac{8\,ka}{3\pi}
=\rho_0 S_d\frac{8 a}{3\pi}\,s\,.}
\]
Therefore the low-frequency added mass (kg) is
\[
\boxed{\,M_{a,\mathrm{LF}}=\frac{X_{\mathrm{rad}}(s)}{s}
=\rho_0 S_d\frac{8 a}{3\pi}
=\frac{8}{3}\rho_0 a^3\,.}
\]
(Equivalently \(M_a=\rho_0 S_d \varepsilon_{\mathrm{LF}}\) with \(\varepsilon_{\mathrm{LF}}=8a/(3\pi)\).)

\subsection{Where the LF expansion breaks down (practical guidance).}
\begin{itemize}
  \item The \((ka)^2/2\) resistance and \(8ka/(3\pi)\) reactance are the first nonzero terms of the Bessel/Hankel expansion and are valid only for \(|ka|\ll1\).  
  \item Conservative rule-of-thumb: \(|ka|\lesssim 0.3\) for errors \(\lesssim 10\%\). For \(ka\sim O(1)\) diffractive and directivity effects dominate; higher-order oscillatory terms from the exact expression become important.  
  \item The reactive (added-mass) term is near-field stored energy and shifts the mechanical resonance \(f_s\); incorrect added-mass factors change low-frequency SPL / resonance predictions.
\end{itemize}

\subsection{Far-field on-axis algebraic simplification (baffled case).}
When the observer is in the far field (Fraunhofer condition; typically \(r\gg a^2/\lambda\)), the spatial integral approximately factors and the on-axis pressure becomes monopole-like with the half-space prefactor (baffle):
\[
p_{\mathrm{piston,on\text{-}axis}}(r,s)\approx \rho_0\,s\,\frac{S_d}{2\pi r}\,v_{\mathrm{avg}}(s)\,e^{-s r/c}.
\]
Hence the far-field baffled-on-axis acoustic impedance is
\[
\boxed{\,Z_{\mathrm{baffled,on\text{-}axis}}^{\mathrm{ac}}(s,r)
=\rho_0\,s\,\frac{S_d}{2\pi r}\,e^{-s r/c}\,.}
\]
(Physically: full-space monopole has denominator \(4\pi r\); the baffle restricts radiation to half-space giving the \(2\pi r\) factor.)\\

\textbf{When it is safe to drop \(e^{-s r/c}\) like in the classic TS model?}
Dropping the exponential is acceptable when:
\begin{itemize}
  \item the model targets low frequencies where the phase delay across the band is negligible, or
  \item the application is sensitivity-level estimation, low-order control design, or steady-state SPL prediction where exact phase is not required, and
  \item \(r\) is fixed and large enough that the delay manifests mainly as an (almost) constant group delay over the modeled band.
\end{itemize}
If causal time-domain accuracy or controller-phase matching is required, keep \(e^{-s r/c}\) or substitute a Pad\'e approximation.

\subsection{Converting acoustic operator to mechanical radiation impedance for motional branch.}
If the mechanical model uses the integrated radiated force, under uniform-face-pressure
\[
Z_{\mathrm{rad}}^{\mathrm{mech}}(s)\approx S_d\,Z_{\mathrm{rad}}^{\mathrm{ac}}(s).
\]
Therefore inserting the far-field algebraic form into the mechanical branch gives (delay dropped if desired)
\[
Z_{\mathrm{rad,mech}}(s)\approx S_d\cdot\frac{\rho_0 s S_d}{2\pi r} = \frac{\rho_0 s S_d^2}{2\pi r},
\]
which shows explicitly the \(S_d^2\) scaling in the mechanical domain when using an acoustic numerator proportional to \(S_d\).

\subsection{Ring-based numerical BEM (full \(ka\), arbitrary \(v(r)\)).}
For accuracy at all \(ka\), for non-uniform velocity, or for finite-baffle/array interactions discretize the radiator into \(N\) rings (areas \(dS_i\)) and use the causal Laplace-domain Green's function \(G(r_i,r_j,\theta)=e^{-kR_{ij}}/R_{ij}\) with
\[
R_{ij}=\sqrt{r_i^2+r_j^2-2r_ir_j\cos\theta}.
\]
The Laplace-domain mutual mechanical impedance between rings \(i,j\) (with \(M\)-point angular quadrature) is
\[
Z_{ij}(s)\approx\frac{s\rho_0}{4\pi}\,dS_i\,dS_j\;\frac{1}{M}\sum_{m=1}^M\frac{e^{-kR_{ij,m}}}{R_{ij,m}},
\]
and
\[
Z_{\mathrm{rad,BEM}}(s)=\sum_{i=1}^N\sum_{j=1}^N Z_{ij}(s),\qquad
M_a(s)=\frac{\Im\{Z_{\mathrm{rad,BEM}}(s)\}}{s}.
\]

\subsection{Implementation and numerical notes}
\begin{itemize}
  \item Treat the weak self-term ( \(i=j\) ) analytically or with refined quadrature; verify convergence by mesh refinement.  
  \item Evaluate Bessel/Y functions for complex arguments with a robust special-function library when using the exact piston form.  
  \item When building rational approximants (Padé, vector fitting) for time-domain simulation compute on the Laplace axis so causality is preserved.  
  \item Explicitly state where you drop \(e^{-s r/c}\) (delay) — mixing delay-free numerators with high-fidelity mechanical \(Z_m^{\mathrm{tot}}(s)\) is permissible but must be justified.
\end{itemize}

\subsection{Key Takeaways}
\begin{itemize}
    \item \textbf{Constant Model:} Best for low frequencies; very simple and fast.
    \item \textbf{Bessel Function Model:} Suitable for high frequencies; moderate complexity.
    \item \textbf{BEM Model:} Accurate across all frequencies; moderate complexity.
\end{itemize}





\newpage
\section{Unified Transfer Function Derivation}

\subsection{Voltage-to-pressure and electrical impedance (closing the loop)}

The electromechanical coupling gives the average cone velocity
\[
v_{\mathrm{avg}}(s)=\frac{B\ell}{Z_m^{\mathrm{tot}}(s)}\,I(s)
=\frac{B\ell}{Z_m^{\mathrm{tot}}(s)}\cdot\frac{V_{\mathrm{in}}(s)}{Z_e(s)},
\]
so the voltage-to-velocity transfer function is
\[
H_v(s)\triangleq\frac{v_{\mathrm{avg}}(s)}{V_{\mathrm{in}}(s)}=\frac{B\ell}{Z_m^{\mathrm{tot}}(s)\,Z_e(s)}.
\]

The acoustic pressure at the listening point is obtained from the appropriate acoustic radiation impedance:
\[
p(s)=Z_{\mathrm{rad}}^{\mathrm{ac}}(s)\;v_{\mathrm{avg}}(s),
\]
hence the voltage-to-pressure transfer function is
\[
H_p(s)\triangleq\frac{p(s)}{V_{\mathrm{in}}(s)}=Z_{\mathrm{rad}}^{\mathrm{ac}}(s)\;H_v(s).
\]

\subsection{Total mechanical impedance}

The total mechanical impedance seen by the motor is
\[
\boxed{
Z_m^{\mathrm{tot}}(s)= sM_{\mathrm{ms}} + R_{\mathrm{ms}} + \frac{1}{sC_{\mathrm{ms}}} + Z_{\mathrm{rad}}(s)
}
\]
where \(Z_{\mathrm{rad}}(s)\) is the \textbf{mechanical radiation impedance} (force per average velocity).  
We decompose \(Z_{\mathrm{rad}}(s)\) into a frequency‑dependent inertial part (added mass) and a residual part:
\[
\boxed{
Z_{\mathrm{rad}}(s) = s\,M_{\mathrm{ma}}(s) + Z_{\mathrm{rad,other}}(s)
}
\]
Here:
\begin{itemize}
  \item \(M_{\mathrm{ma}}(s)\) = frequency‑dependent added mass obtained from the Kirchhoff (Rayleigh) integral solution (method‑dependent).
  \item \(Z_{\mathrm{rad,other}}(s)\) = remaining mechanical impedance, typically the radiation resistance \(R_{\mathrm{rad}}(s)\) plus enclosure loading, horn effects, or other structural contributions.
\end{itemize}

Thus the total mechanical impedance becomes
\[
Z_m^{\mathrm{tot}}(s)= s\bigl(M_{\mathrm{ms}}+M_{\mathrm{ma}}(s)\bigr) + R_{\mathrm{ms}} + \frac{1}{sC_{\mathrm{ms}}} + Z_{\mathrm{rad,other}}(s).
\]

\subsection{Electrical impedance}

The electrical impedance seen at the voice‑coil terminals follows from the reflected motional term:
\[
\boxed{
Z_e(s) = R_e + sL_e + \frac{(B\ell)^2}{Z_m^{\mathrm{tot}}(s)}.
}
\]

Substituting the expression for \(Z_m^{\mathrm{tot}}(s)\) gives the explicit form
\[
\boxed{
Z_e(s) = R_e + sL_e +
\frac{(B\ell)^2}{
s\bigl(M_{\mathrm{ms}} + M_{\mathrm{ma}}(s)\bigr)
+ R_{\mathrm{ms}} + \dfrac{1}{sC_{\mathrm{ms}}}
+ Z_{\mathrm{rad,other}}(s)
}
}.
\]

\subsection{Voltage-to-pressure transfer function (maximally correct causal form)}

Using the far‑field algebraic approximation for a baffled piston, including the propagation delay and general solid angle \(\Omega\), we obtain
\[
\boxed{
H_p(s) = \frac{p(s)}{V_{\mathrm{in}}(s)}
= \frac{ s \;\cdot\; \dfrac{B\ell \,\rho_0 S_d}{\Omega \, r} \; e^{-s r/c} }
{ (B\ell)^2 \;+\; (R_e + s L_e) \left[ s\bigl(M_{\mathrm{ms}} + M_{\mathrm{ma}}(s)\bigr) + R_{\mathrm{ms}} + \dfrac{1}{s C_{\mathrm{ms}}} + Z_{\mathrm{rad,other}}(s) \right] }
}.
\]

Where:
\begin{enumerate}
  \item \(M_{\mathrm{ma}}(s)\) = frequency‑dependent added mass (reactive part of the radiation load, from Kirchhoff/Rayleigh).
  \item \(Z_{\mathrm{rad,other}}(s)\) = remaining mechanical radiation impedance (e.g., radiation resistance \(R_{\mathrm{rad}}(s)\) plus enclosure, horn, or other effects).
  \item \(\Omega\) = solid angle of radiation (\(2\pi\) for half‑space/infinite baffle, \(4\pi\) for free field, \(\pi\) for a corner, etc.).
  \item \(e^{-s r/c}\) = causal propagation delay from driver to listener at distance \(r\).
\end{enumerate}

\subsection{Relation to classical T/S model}
If we neglect the propagation delay (or absorb it into a constant phase), take \(\Omega = 2\pi\) (infinite baffle), drop \(Z_{\mathrm{rad,other}}(s)\), and approximate \(M_{\mathrm{ma}}(s)\) by its low‑frequency constant value \(M_{\mathrm{ma}} = \frac{8}{3}\rho_0 a^3\), the expressions reduce to the classic Thiele–Small voltage‑to‑pressure transfer function and electrical impedance.

\subsection{Practical modeling notes}
\begin{itemize}
  \item Use the full Rayleigh integral when per‑element accuracy or non‑uniform face velocity is important (evaluate numerically via ring‑BEM or panel BEM).  
  \item Use the monopole form when the radiator is electrically/physically small relative to wavelengths of interest.  
  \item Use the baffled far‑field algebraic factor for simple on‑axis sensitivity and control‑oriented models; include Padé approximations for the causal delay if you need a rational model.  
  \item When converting between acoustic and mechanical impedances, be careful: \(Z_{\mathrm{rad}}^{\mathrm{mech}}(s)\) (N·s/m) relates to \(Z_{\mathrm{rad}}^{\mathrm{ac}}(s)\) (Pa·s/m) through surface‑area integration and is only simply \(S_d Z_{\mathrm{rad}}^{\mathrm{ac}}(s)\) under the uniform‑pressure approximation.
\end{itemize}

\subsection*{Key physical relationships}
\begin{align*}
\text{Acceleration to Pressure:} &\quad p(s)=\underbrace{\frac{\rho_0 S_d}{2\pi r}}_{K_{\text{rad}}}\,s^2 x(s), \\
\text{Velocity to Pressure:} &\quad p(s)=\underbrace{\frac{\rho_0 S_d s}{2\pi r}}_{Z_{\mathrm{rad}}^{\mathrm{ac}}(s)}\,v(s), \\
\text{Position to Pressure:} &\quad p(s)=s^2 K_{\text{rad}} x(s).
\end{align*}

\newpage
\section{Region Overview}

\subsection{Frequency-Domain Physics and Design Implications}
\begin{center}
\small
\begin{tabular}{|c|c|c|l|}
\hline
\textbf{Frequency Range} & \textbf{Dominant Physics} & \textbf{Pressure relation (dominant)} & \textbf{Design Focus} \\
\hline
$f < f_s$ & Suspension compliance dominates & $p\propto s^2 x$ (radiation small) & \begin{tabular}[c]{@{}l@{}}- Excursion limits\\- $X_{\text{max}}$ optimisation\\- High-compliance suspension\end{tabular} \\
\hline
$f_s < f < f_e$ & Mass control (inertia) & $p\propto s v$ & \begin{tabular}[c]{@{}l@{}}- Sensitivity optimisation\\- Motor strength $(B\ell)$\\- Low $M_{\text{ms}}$\end{tabular} \\
\hline
$f > f_e$ & Voice coil inductance & electrical rolloff affects drive & \begin{tabular}[c]{@{}l@{}}- Shorting rings\\- Copper caps\\- Crossover compensation\end{tabular} \\
\hline
$f > f_{\text{rad}}$ & Radiation directivity / lobing & $p$ shows strong $ka$ dependence & \begin{tabular}[c]{@{}l@{}}- Piston size optimisation\\- Dome vs cone tweeters\\- Phase plugs\end{tabular} \\
\hline
\end{tabular}
\end{center}

where
\begin{align*}
f_s &= \frac{1}{2\pi\sqrt{(M_{\mathrm{ms}} + M_a)C_{\mathrm{ms}}}} & \text{(Resonance frequency)} \\
f_e &= \frac{R_e}{2\pi L_e} & \text{(Inductive rolloff, rough)} \\
f_{\text{rad}} &= \frac{c}{2\pi a} & \text{(frequency at which $ka\sim1$)} \\
a &= \sqrt{S_d/\pi} & \text{(Equivalent piston radius)}
\end{align*}

\subsection{Speakers in Series (isobaric) and Parallel}
\textbf{Assumption: $r\ll\lambda$ and isobaric cavity mass negligible.}
\small
\begin{center}
\begin{tabular}{lccccc}
\toprule
\textbf{Configuration} & $\mathbf{S}$ & $\mathbf{C}$ & $\mathbf{V_{\mathrm{as}}}$ & $\mathbf{P_{\max}}$ & $\Delta L_p$ \\
 & (Radiating area) & (Compliance) & (Eq. volume) & (Power handling) & (vs. single) \\
\midrule
Single driver & 
$S$ & 
$C$ & 
$V_{\mathrm{as}}$ & 
$P$ & 
$0\,\mathrm{dB}$ \\

Isobaric/Series pair & 
$S$ & 
$\tfrac{C}{2}$ & 
$\tfrac{V_{\mathrm{as}}}{2}$ & 
$2P$ & 
$\begin{cases}
0\,\mathrm{dB}, & P\ (\text{same total})\\
+3\,\mathrm{dB}, & 2P\ (\text{doubled})
\end{cases}$ \\

Two drivers in parallel & 
$2S$ & 
$2C$ & 
$2V_{\mathrm{as}}$ & 
$2P$ & 
$\begin{cases}
+3\,\mathrm{dB}, & P\ (\text{same total})\\
+6\,\mathrm{dB}, & 2P\ (\text{doubled})
\end{cases}$\\
\bottomrule
\end{tabular}

\medskip
\textbf{Table:} Comparison of a single driver, an isobaric pair, and two drivers in parallel (coherent; otherwise -3 dB when incoherent)
\end{center}

\newpage
\subsection{Source Addition in Spherical Speaker Arrays}

\subsubsection{Principles at the Listening Point}
\begin{itemize}
  \item \textbf{Coherent sources:} sources with fixed phase and frequency relationships; sound pressures add linearly: $p_{\mathrm{total}}=N\cdot p_{\mathrm{single}}$. Level increases by $20\log_{10}N$.  
  \item \textbf{Incoherent sources:} intensities add, pressures combine statistically as root-sum-square: $p_{\mathrm{total}}=\sqrt{N}\,p_{\mathrm{single}}$. Level increases by $10\log_{10}N$.
\end{itemize}

\subsubsection{Arbitrary Speaker Scaling by Solid Angle}
For a sound source with known reference SPL $L_{p,\mathrm{ref}}$ at $P_{\mathrm{ref}}=1\,\mathrm{W}$, $r_{\mathrm{ref}}=1\,$m and reference solid angle $\Omega_0=4\pi$ the general scaling is
\[
L_p(r,P,\Omega)=L_{p,\mathrm{ref}}+10\log_{10}\left(\frac{P}{P_0}\right)-20\log_{10}\left(\frac{r}{r_0}\right)+10\log_{10}\left(\frac{\Omega_0}{\Omega}\right).
\]

\begin{center}  % centers the table
\begin{tabular}{lcc}
\toprule
Geometry & Solid Angle $\Omega$ & Level Gain [dB] \\
\midrule
Spherical wave (free field)       & $4\pi$          & $0\,\mathrm{dB}$  \\
Half-space (baffle)               & $2\pi$          & $+3.01\,\mathrm{dB}$ \\
Quarter-space (corner)            & $\pi$           & $+6.02\,\mathrm{dB}$ \\
Eighth-space (floor \& corner)   & $\frac{\pi}{2}$ & $+9.03\,\mathrm{dB}$ \\
\bottomrule
\end{tabular}
\end{center}


% Add References
\newpage
\subsection*{Short reference list}
Lord Rayleigh, \emph{The Theory of Sound}, Vol. I (Macmillan, 1878).\\
Leo L. Beranek, \emph{Acoustics} (McGraw-Hill, 1954).\\
Robert H. Small, "Direct-Radiator Loudspeaker System Analysis," \emph{J. Audio Eng. Soc.}, vol. 20, no. 10, 1972.



\end{document}
